\section{Discussion}
\label{sec:discussion}

\subsection{Blaming and Defending Dissent Over Time}
\label{subsec:discussion-temporal-dissent}
The temporal results show that minority voice did not always decline as the discussion progressed. In AITA, blaming dissent generally decreased, while defending dissent increased in the all-agreement and low-agreement discussions. In high-agreement discussions, defending dissent declined at first, but then increased and ended at a level close to where it started.

Agreement level provides further context for these trajectories. Changes in blaming and defending shares were larger in low-agreement AITA discussions, while high-agreement discussions showed smaller absolute changes. The certainty results provide a further observation. Both types of dissent were expressed with greater certainty toward the end of the discussion. Thus, later minority comments were expressed with greater certainty, even when their share of the discussion did not increase.

The community comparison sets a boundary on these findings. Across all agreement levels combined, AITAH showed the opposite directional pattern: blaming share increased while defending share decreased. This suggests that the temporal development of moral dissent may differ across communities. Taken together, these observations suggest that minority participation should be understood in relation to both the direction of dissent and the community in which it occurs.

\subsection{Interpreting Responses to Minority Judgments}
\label{subsec:discussion-minority-responses}

\textbf{[TODO: develop the discussion along the following outline.]}
\begin{enumerate}
    \item \textbf{Beyond reply volume: what responses do minority judgments receive?}
    \begin{itemize}
        \item Connect the score and reply-volume findings to the response-type analysis.
        \item Explain how distinguishing support, challenge, information-seeking, and
    other responses helps interpret engagement with dissent: more interaction
    does not necessarily mean greater agreement.
        \item Keep the composition of observed replies distinct from the probability that a comment receives
    a reply or a challenge.
    \end{itemize}

    \item \textbf{Responder roles and the direction of dissent.}
    \begin{itemize}
        \item Connect the interpretation to blaming and defending dissent, and consider
    how the OP's role as the person being judged may differ from that of other
    participants.
        \item If OP replies are included as examples, present them separately from the
    annotation statistics, which exclude OP participation.
    \end{itemize}
\end{enumerate}
